% Transcription — fonds Alexandre Grothendieck, shelfmark 154, batch 4 (pages 61-80).
% Established from the facsimile archives/batches/154/batch-04.pdf, cut from a
% copy of 154.pdf fetched through the project's own relay
% (https://grothendieck-relay.onrender.com/source/154.pdf); PDF page k =
% archivists' page 60+k, no cover sheet. Per the skill transcribe-grothendieck.
% The folder is 175 pages in nine batches (1-20, ..., 161-175), transcribed in
% parallel; batches 2, 3 and 5 were read for notation and joins (p. 61
% continues the « 4.1. » alinea at the end of batch 3's p. 60).
% Notation follows batches 1-3, 5 and the transcription of folder 80.
%
% Pass: Opus 5.5 (claude-opus-5-5), 2026-09-26 — first pass, unchecked against the pages by a human.
%
% Dating and title. \dating is the folder's own date, copied verbatim from
% src/content/catalogue.ts; the folder belongs to the inventory group
% « Espaces stratifiés » (150-151), but since it has a date of its own the
% group's range is not added. Title from the catalogue, trailing full stop
% dropped; « [Système de pseudo-droites] » is bracketed there because the
% archivists supplied it.
%
% Paper. Fan-fold computer-listing paper written in landscape, two columns to
% a leaf; pp. 61 and 72 are pencil notes, p. 72 written with the listing
% upside down (its JCL banner is paper, not a date). Not described further.
%
% Pages skipped: none. Pages summarised: none. Page given only as a figure
% described in a French \note, with its caption: 80. Pages largely pencil
% calculation or sketch: 61, 72. All others (62-71, 73-79) transcribed from
% his hand, fast drafting with many \ill{} and \uncertain{}.
%
% His pagination and dates. No page number of his is visible. One date in
% his hand: « 3.1.84 » at the head of p. 64 (same day as batch 3's p. 57).
% Numbered runs: the two kinds of fundamental circuits a) b) on p. 63; the
% « huit façons » of gluing Delta_D to Delta_D' (p. 64) and the resulting
% list A) 1) 2) (p. 68) and B) 1) 2) (p. 69), referred to as A1, A2, B1, B2
% on pp. 69-70; the two mistakes 1) 2) of the previous two weeks on p. 67;
% « Description combinatoire » 1) 2) 3) on pp. 75-76. All recorded as
% written; no numbering of his is continued across the batch boundary.
%
% What the batch is about.
%   61-63  (continuing 4.1 of batch 3) Parallel transport of orientations
%          between faces F_D, F_D'; the characters chi_*, chi_s, chi_a of
%          pi_1(X, F_D) with chi_* = chi_s chi_a, and the two further
%          characters chi_X (orientation of the surface) and chi_c
%          (« combinatoire », parity of paths); values on the fundamental
%          circuits (u(l_{D_i}) = translation by n-1, half-circuits round a
%          point), whence pi_1(X, F_D) -> Aut(F_D) surjective (« ouf ! »);
%          chi_c trivial by the parity of sing_D, whence chi_X = chi_*.
%   64-70  3.1.84: the eight « standard » ways of gluing Delta_D to Delta_D'
%          along a transition arc (a, a' against b, b'), reduced to four;
%          shrinking to Delta*_D so that the glued surface maps to X;
%          alpha = a' gives an immersion X -> X; the transition isomorphism
%          on orientations is opposite to parallel transport; his admission
%          that for two weeks he had « mélangé les pinceaux deux fois »;
%          the list A1, A2 (natural map X -> X) and B1, B2 (none), checked
%          against the circuits round a supersingular position: only B2
%          fails.
%   71-76  Gluing with stable relative positions only: symmetric figure,
%          couples (a, b), (a', b'); case « de prolongation des rives » with
%          parallel transport (immersion, K, L = phi^{-1}(X_1)) or its
%          opposite (the right supersingular circuits); a second polygonal
%          structure from the disjoint specialisation arcs on D, n-gons and
%          a Möbius band; the combinatorial description 1)-3) (graph Gamma of
%          stable positions, polygonal structure, orientation transport) and
%          the variant with folds along the edges of K. P. 72 is a pencil
%          aside: sigma_0 ... sigma_{N-1}, N = 4n, rho_i = sigma_{i-1}sigma_i.
%   77-80  The case of a system of straight lines in a real projective plane
%          X: the dual arrangement in X^v of the lines joining the points
%          delta_i, faces F of the dual system and their edges a_alpha as
%          specialisations D -> D_alpha; pivoting D about s_alpha to follow
%          a_alpha, antiscompatible orientations; the absence of a natural
%          map X -> X compared with that of X^v -> X; p. 80 drawings of
%          half-circuits round points of order 1, 2 or 0.
%
% Notation fixed once for the batch, following batches 3 and 5:
% \mathcal{X} for his script X (the glued surface), X for the ambient
% surface; \Delta_D, \Delta^{*}_D (the shrunken disc); a, a', b, b' for
% transition arcs and their antipodes; \underline{\omega} for his underlined
% omega; \chi_{\ast}, \chi_s, \chi_a, \chi_{\mathcal{X}}, \chi_c and
% \varepsilon_s, \varepsilon_a as in batch 3's p. 60; \mathbb{D}_2;
% \underline{\Sigma}, \check{\underline{\Sigma}} for the dual system;
% \delta_i, \Delta_\alpha, s_\alpha, a_\alpha.
%
% Readings to check first: the indices of the characters on p. 62 (* versus
% X, a versus o); the end of p. 63; pp. 64-65 (the triples, the cancelled
% head of p. 65); p. 69 (« bâton-dans-les-roues »); pp. 74-75 (fast, many
% doubtful words); p. 79 (« antiscompatible »).
\documentclass[11pt,a4paper]{article}
% Le préambule commun à toutes les transcriptions du fonds Grothendieck.
%
% Il tient en une idée : **l'incertitude se compose, elle ne se résout pas.**
% Une transcription qui lisse un mot douteux a détruit la seule chose pour
% laquelle on la faisait — un lecteur doit pouvoir savoir, à chaque endroit,
% ce qui était sur le feuillet et ce qui a été ajouté. Les six macros qui
% suivent sont donc l'appareil critique, et non de la décoration.
%
% Les mêmes commandes sont reconnues par `scripts/render.mjs`, qui produit la
% vue de lecture du site. Une macro ajoutée ici doit l'être aussi là-bas,
% faute de quoi le rendu échouera bruyamment — ce qui est voulu : un rendu
% silencieusement faux est pire qu'un rendu absent.

\ProvidesPackage{grothendieck}[2026/08/08 Transcriptions du fonds Grothendieck]

\usepackage{amsmath,amssymb,amsthm}
\usepackage{mathrsfs}
\usepackage{stmaryrd}
% Les diagrammes commutatifs sont partout dans ce fonds : c'est la langue
% même de la géométrie algébrique de Grothendieck, et une transcription qui
% les rendrait en prose ne transcrirait rien.
\usepackage{tikz-cd}
% Français par défaut — c'est la langue des feuillets — mais l'anglais est
% chargé aussi : la traduction et le résumé sont des documents anglais, et la
% typographie française (espaces avant les deux-points) y serait une faute.
% Ces documents basculent par \selectlanguage{english} en tête de corps.
\usepackage[english,french]{babel}
\usepackage{fontspec}
\usepackage{geometry}
\geometry{a4paper,margin=25mm,marginparwidth=28mm}
\usepackage{marginnote}
\usepackage{xcolor}
% ulem, not soul. `soul`'s \st and \ul are built on a character-by-character
% scan that XeLaTeX's font machinery breaks: the document fails with an
% undefined control sequence rather than degrading. ulem does the same two jobs
% and is Unicode-engine safe. `normalem` stops it hijacking \emph, which the
% transcriptions use for Grothendieck's own emphasis.
\usepackage[normalem]{ulem}
\usepackage{hyperref}
\hypersetup{colorlinks=true,linkcolor=black,urlcolor=black,pdfborder={0 0 0}}

% --- Métadonnées du lot -----------------------------------------------------
% Reprises telles quelles par le script de rendu, qui les affiche en tête de
% la vue de lecture. Elles fixent ce que le fichier prétend couvrir : sans
% elles, une transcription flotte sans point d'attache dans un fonds de seize
% mille feuillets.

\newcommand{\folder}[1]{\gdef\grothFolder{#1}}
\newcommand{\batch}[1]{\gdef\grothBatch{#1}}
\newcommand{\leaves}[2]{\gdef\grothFirst{#1}\gdef\grothLast{#2}}
\newcommand{\foldertitle}[1]{\gdef\grothFoldertitle{#1}}
\gdef\grothFolder{?}\gdef\grothBatch{?}\gdef\grothFirst{?}\gdef\grothLast{?}\gdef\grothFoldertitle{}

% --- Le résumé --------------------------------------------------------------
%
% Il ouvre la lecture modernisée, sous le titre « Résumé », et il est composé
% en petit. Ce n'est pas un ornement : c'est le seul passage du document qui
% ne soit pas des mathématiques, et un lecteur doit pouvoir le reconnaître
% avant de l'avoir lu — puis le sauter s'il connaît déjà le sujet, ou s'y
% arrêter s'il ne le connaît pas. Le filet qui le suit marque la reprise du
% corps.

\newenvironment{resume}
  {\small\setlength{\parskip}{0.45em}}
  {\par\medskip\hrule\medskip}

% --- Mots-clés ---------------------------------------------------------------
%
% En anglais, en clôture du résumé de la lecture modernisée : le vocabulaire
% moderne sous lequel chercher ce que les feuillets construisent. Le manifeste
% du site (`npm run manifest`) les extrait pour étiqueter la cote entière —
% cette ligne est la source unique des tags, il n'y a pas de fichier à côté.
\newcommand{\keywords}[1]{\par\smallskip\noindent
  {\footnotesize\textsc{Keywords} --- \textit{#1}}\par}

% --- L'appareil critique ----------------------------------------------------

% Le numéro de feuillet, en marge. C'est l'ancre qui permet de régler un
% désaccord sur une formule en regardant *un* feuillet plutôt qu'une chemise
% de six cents. Le site s'en sert aussi pour tourner les pages du fac-similé
% quand on fait défiler la transcription.
\newcommand{\leaf}[1]{%
  \marginnote{\footnotesize\color{gray!70}\textsf{#1}}%
  \label{leaf:#1}\ignorespaces}

% Illisible : on ne devine pas. Un mot inventé qui se lit comme les autres est
% le pire résultat possible.
\newcommand{\ill}{\textcolor{red!60!black}{[\,\ldots\,]}}

% Lecture douteuse : le mot est donné, et signalé comme incertain.
\newcommand{\uncertain}[1]{\uline{#1}}

% Ajout de l'éditeur — une lettre manquante, un mot sous-entendu.
\newcommand{\add}[1]{\textcolor{blue!50!black}{[#1]}}

% Biffé par Grothendieck lui-même. Ce qu'il a rayé fait partie du document :
% une rature dit souvent où la pensée a changé de direction.
\newcommand{\struck}[1]{\sout{#1}}

% Note du transcripteur, sur l'état matériel du feuillet.
\newcommand{\note}[1]{\par\smallskip{\small\color{gray!60!black}\textit{[#1]}}\par\smallskip}

% Note marginale de Grothendieck, distincte du corps du texte.
\newcommand{\marginal}[1]{\par\smallskip\noindent\hspace{1em}%
  {\small\color{gray!60!black}#1}\par\smallskip}

% --- Titre ------------------------------------------------------------------

\AtBeginDocument{%
  \begin{center}
    {\large\bfseries Fonds Alexandre Grothendieck --- cote n\textsuperscript{o}~\grothFolder}\\[2pt]
    {\itshape\grothFoldertitle}\\[4pt]
    {\small feuillets \grothFirst--\grothLast\ (lot \grothBatch)}\\[2pt]
    {\footnotesize\color{gray!60!black}Transcription établie d'après le fac-similé numérisé
     par l'Université de Montpellier.\\
     Les lectures incertaines sont soulignées, les illisibles notées [\,\ldots\,],
     les ajouts entre crochets.}
  \end{center}
  \vspace{1em}}

\endinput


\folder{154}
\batch{4}
\pages{61}{80}
\dating{1983-1984}
\watermark{Édition de démonstration}
\foldertitle{[Système de pseudo-droites] : notes manuscrites (1983-1984, s.d.)}

\begin{document}

\page{61}
\[
F = F_D \longrightarrow F_{D'} = F'
\]
\[
\begin{array}{ccc}
\omega(F_D) & \xrightarrow[\text{cond.\ via les } \operatorname{Pol}(D)]{\alpha \ \text{transport parallèle}} & \omega(F_{D'}) \\
\wr & & \wr \\
\omega(F) & \xrightarrow[\text{l'orientation sur } \mathcal{X}]{\beta \ \text{transport de}} & \omega(F') \\
\wr & & \wr \\
\omega(D) & \xrightarrow{\alpha' \ \text{transport parallèle}} & \omega(D')
\end{array}
\]
\note{notes au crayon, écrites par-dessus un listing d'ordinateur ; elles reprennent, pour un chemin de $F = F_D$ à $F' = F_{D'}$, le transport des orientations de l'alinéa « 4.1. » de la page 60 (lot 3). Le « cond.\ » sous la première flèche est d'une lecture \uncertain{douteuse} ; dans $\omega(F)$ de la deuxième ligne, le $F$ est écrit sur une autre lettre}

\[
\alpha = \alpha' \qquad \beta = \alpha\alpha' \ \text{ssi nb de } \struck{\text{points}}\ \text{\uncertain{sing.}}\ N \text{ est pair}
\]

\emph{ex} $F = F_{D'}$, $\alpha = \alpha' = -1$ \quad alors $\beta = -1$ ssi tel $N$ pair

donc
\note{la page s'arrête sur « donc ». « $N$ » est ajouté au-dessus de la ligne ; le mot biffé qui le précède se lit mal. Figures au crayon, à droite : un cercle pincé en huit, une flèche en boucle le long de son bord ; au-dessous, deux droites horizontales raccordées par un palier en S, traversées de trois verticales, deux points noirs aux croisements}

\page{62}
et tout chemin $c$ définit des iso.\ d'\uncertain{une} \ill{} équivariants
\[
\begin{array}{ccc}
\underline{\omega}(F) \simeq \underline{\omega}(D) & \xrightarrow{\ \underline{\omega}_{\ast}(c)\ } & \underline{\omega}(F') \simeq \underline{\omega}(D') \\[4pt]
\varepsilon_s(F) = \varepsilon_s(D) & \longrightarrow & \varepsilon_s(F') \simeq \varepsilon_s(D') \\[4pt]
\varepsilon_a(F) \simeq \varepsilon_a(D) & \longrightarrow & \varepsilon_a(F') \simeq \varepsilon_a(D')
\end{array}
\]
\note{la page continue la fin de la page 60 (lot 3), où sont introduits $\chi_{\ast}$ et les $\varepsilon_s(F)$, $\varepsilon_a(F)$ (« sommets », « arêtes ») ; nous gardons ces notations. Les lettres rendues $\varepsilon$ sont des majuscules en forme d'epsilon ou de sigma ; l'indice rendu $a$ est un petit rond, et l'indice rendu $\ast$ (ici et dans $\underline{\omega}_{\ast}$, $\chi_{\ast}$) une petite croix, lue comme au lot 3 --- lectures \uncertain{probables}. Le « $c$ » de « chemin $c$ » est ajouté au-dessus de la ligne}

Les iso.\ ne sont pas indépendants, \uncertain{de quoi} pour en deux déterminer le troisième, par la relation
\[
\underline{\omega}(F) \simeq \varepsilon_s(F) \wedge \varepsilon_a(F)
\]
et les deux relations similaires, équivalentes. Au niveau des caractères correspondants sur $\pi_1(\mathcal{X}, F)$, cela donne
\[
\boxed{\chi_{\ast}(F) = \chi_s(F)\,\chi_a(F)}
\]
\[
\begin{array}{c}
\| \\
\chi_{\ast}(D) = \chi_s(D)\,\chi_a(D) \\
\| \\
\underline{\omega}(D)
\end{array}
\]
et on a \ill{} plus symétriquement que
\[
\chi_{\ast}(F)\,\rho_s(F)\,\chi_a(F) = 1.
\]
\note{dans la dernière formule, le deuxième facteur est écrit $\rho_s$ plutôt que $\chi_s$ ; nous transcrivons tel quel}

En plus de ces trois caractères, il y en a deux autres, $\chi_{\mathcal{X}} = {}$caractère d'orientation de $\mathcal{X}$ (\uncertain{qui diffère} de $X$ par la structure cellulaire), et $\chi_c$ (« combinatoire »), qui est défini par la parité\supplied{ des longueurs}. Le \struck{système} \add{\uncertain{sous-système}} local correspondant \struck{\ill{}} défini par les $(\underline{\omega}(F))$, mais ceci était \ill{} déduit des transitions \emph{\uncertain{opposées}} à celles données par \struck{\ill{}} le transport parallèle sur $X$ défini par les $\underline{\omega}_{\ast}(c)$ --- donc on trouve
\[
\underline{\omega}_{\mathcal{X}}(c) = \operatorname{par}(c)\,\underline{\omega}_{\ast}(c)
\]
où $\operatorname{par}(c)$ désigne la parité de $c$. \struck{Le système} Le système local associé à $\chi_c$ \struck{formé de} a comme fibres les ens.\ $\lbrace \pm 1 \rbrace$, mais les morphismes de transition donnés par les $\operatorname{par}(c)$. On trouve donc
\[
\boxed{\chi_{\mathcal{X}} = \chi_a\,\chi_c}
\]
\note{dans la formule encadrée, le deuxième indice est le même petit rond que plus haut ; la formule précédente donnerait plutôt $\chi_{\ast}\chi_c$, comme à la dernière ligne de la page. Nous transcrivons ce que porte la page}
\note{« des longueurs » : ajout interlinéaire d'une lecture \uncertain{douteuse}}

Compte tenu de ces \ill{}, on dispose de trois caractères fondamentaux, \add{\uncertain{privilégiés}} je choisis
\[
\chi_{\ast}, \quad \chi_s, \quad \chi_c
\]
\ill{} les deux autres \ill{}
\[
\chi_a = \chi_{\ast} \chi_s \quad \text{et surtout} \quad \chi_{\mathcal{X}} = \chi_{\ast} \chi_c.
\]
\marginal{On a $\chi_c = 1$, cf.\ verso !}
\note{note écrite dans la marge de droite et séparée du texte par un long trait oblique ; le « verso » est la page 63}

Il y a lieu d'expliciter les valeurs des caractères précédents, à \uncertain{fortiori}, ainsi que celles de la représentation
\[
u : \pi_1(\mathcal{X}, F_D) \longrightarrow \operatorname{Aut}(\operatorname{Pol}(D)),
\]

\marginal{\uncertain{Attention}, \ill{} \ill{} \emph{\uncertain{fixé}}, dans \ill{} il faut \ill{} \ill{} \emph{\ill{}} \ill{} déterminant \ill{} \ill{} (qui \ill{} \ill{} \ill{} \ill{} \ill{}}
\note{note écrite en biais au bas de la colonne de gauche, reliée par une flèche à la formule $\chi_{\ast}(F)\rho_s(F)\chi_a(F) = 1$ ; elle se lit mal}

\page{63}
Pour les circuits fondamentaux sur $\mathcal{X}$, qui \uncertain{sont} de deux espèces
\note{« fondamentaux » : le mot est écrit dans la marge au-dessus de « a) », avec un trait de renvoi}

a) Circuit autour d'une \struck{\ill{}} \add{\uncertain{composante}} exceptionnelle de $K$ (ou le franchissant une fois), soit $l_{D_i}$ le circuit correspondant à la ps.\ droite $D_i \in \Sigma$. On trouve
\[
u(l_{\vec{D}_i}) = \text{translation par } n-1
\]
(dans le sens de l'orientation de la position rel.\ $D$ voisine de $D_i$), valeurs des \uncertain{caractères} fondamentaux
\[
\chi_{\mathcal{X}}(l_{D_i}) = 1, \quad \chi_s(l_{D_i}) = (-1)^{n-1}
\]
\struck{$\chi_c(l_{D_i}) = 1$}
\note{dans la première valeur, un signe est surchargé en $1$ ; dans la deuxième, un « $=$ » est surchargé et un premier second membre biffé devant $(-1)^{n-1}$ ; la troisième égalité, donnée à la suite, est biffée. L'indice de $\chi$ dans la première est un $\mathcal{X}$ \uncertain{cursif}}

b) Demi-circuit autour d'un pt de $X$ (qui peut être soit dans une face, soit dans une arête, soit \add{être} un sommet). Le \struck{point} fait \ill{} de ce parcours dans $\operatorname{Pol}(D)$ \uncertain{induit} \ill{} par $u(l)$, qui est de signe $-1$, donc on trouve toutes les symétries, donc
\[
u_D : \pi_1(\mathcal{X}, F_D) \longrightarrow \operatorname{Aut}(F_D)
\]
est \emph{surjectif} ouf !
\note{l'alinéa b) est marqué d'un double trait vertical dans la marge de gauche ; « surjectif » est souligné deux fois}

* Il semble que le caractère $\chi_c$ soit \emph{trivial} ! En effet, si pour une p.r.\ $D$ donnée, on désigne par $\operatorname{sing}_D$ (« singularités de $D$ ») le nb de sommets de $\Sigma$ qu'elle contient, je dis que pour deux voisines $D$ et $D'$, on a
\[
\text{poids } c \equiv \operatorname{sing}_{D'} - \operatorname{sing}_D \pmod 2
\]
(donc poids $c$ \emph{pair}, si $D = D'$ !) \ill{} \ill{} \ill{} \ill{} si on \ill{} $c$ \ill{} de longueur $1$, auquel cas la formule est triviale (on \uncertain{se ramène} au cas où la générisation « balaye » au-dessus \struck{d'une droite} \add{d'un sommet}, ou d'une \struck{droite} $D_i \in \Sigma$, le poids est $2$). Donc finalement, on trouve
\[
\boxed{\chi_{\mathcal{X}} = \chi_{\ast}}
\]
il n'y a pas d'autre caractère \ill{} que ceux provenant de $\mathbb{D}_2$ !
\note{l'alinéa est marqué d'un astérisque ; la congruence porte « (2) » à droite, lu comme « modulo $2$ ». « pair » est souligné ; « d'un sommet » est écrit après « d'une droite » biffé. Les mots avant « si on » se lisent mal}

\page{64}
\emph{3.1.84}
\note{date de sa main, en tête de la page, au-dessus des figures}

\note{figures : deux grands disques côte à côte. À gauche, $\Delta_{D'}$, au bord vert foncé (légende $D'$) : deux pseudo-droites orange qui se croisent en $s$ juste au-dessous du bord, en haut, où un segment épais noir marque l'arc $a$ et un arc pointillé l'arc $a'$ ; un petit cercle vert clair y figure le recollement ; une croix orange (point double) près du bord. À droite, $\Delta_D$, au bord vert clair (légende $D$, $\Delta_D$) : deux pseudo-droites orange, un segment épais noir en bas (l'arc $b$, près de $s$) et une rangée de points noirs en haut (l'arc $b'$), une croix orange}

\ill{} huit façons « standard » de recoller $\Delta_D$ à $\Delta_{D'}$, suivant \struck{\ill{}} un isom.\ (combinatoire) d'un des deux arcs de transition $a, a'$ sur $\Delta_{D'}$, avec l'un des deux arcs privilégiés $b$ (de transition) et $b'$ (son antipodique) sur $\Delta_D$, deux choix pour chaque choix d'un couple dans $\lbrace a, a' \rbrace \times \lbrace b, b' \rbrace$. Donc huit façons, \add{\uncertain{cellulairement}} \uncertain{a priori}, de construire une surface : toutes ont le graphe $\Gamma_{\underline{\Sigma}} = \Gamma$, mais (sauf pour les polygones d'un couloir \uncertain{rouge}, \uncertain{peut-être}) il n'y a que quatre cas différents de surfaces, \uncertain{on} sait qu'on obtient par un choix $(\alpha, \beta, \rho)$ \ill{} recollés et $(\alpha', \beta', \sigma')$ (où $\alpha', \beta', \sigma'$ déduits en changeant \ill{} \ill{} les antipodismes) sont can.\ isom.\ (ils correspondent à la \uncertain{même} circularisation de $\Gamma$, et les \uncertain{mêmes} isomorphismes de transition par les or.\ loc.\ correspondantes).
\note{« a priori » est écrit au-dessus de « façons », « cellulairement » au-dessous. Les lettres du troisième terme des triplets, lues $\rho$ et $\sigma'$, sont \uncertain{peu nettes}}

On peut les supposer, quitte à \uncertain{changer} les \uncertain{notations} \ill{} \ill{} et d'autre, que $\beta = b$ --- ce qui revient à faire un choix de la \uncertain{manière} dont s'apparaît \struck{\ill{}} au voisinage des lieux de transition choisis, sur les surfaces recollées, bien qu'il soit \uncertain{associé} \uncertain{en} antipodes, ce qui a l'air un peu idiot !)

\struck{Reste les quatre choix (en fait, les huit) qui}
\note{les deux dernières lignes de la colonne de droite sont barrées de deux traits obliques ; la suite, en tête de la page 65, est elle-même annulée}

\page{65}
\struck{\ill{} \ill{} qu'elles \ill{} \ill{} une application canonique \ill{} \ill{} \ill{} surface recollée dans $X$ ! Si on recolle $b$ à $a'$,} \add{\struck{Non, c'est vrai pourvu \ill{}}} \struck{\ill{} que, par ce choix, on choisisse l'identification}
\note{les cinq premières lignes de la colonne de gauche sont barrées de trois traits obliques ; « Non, c'est vrai pourvu \ldots{} » est écrit au-dessus de « Si on recolle $b$ à $a'$ », lui-même biffé d'un trait}

Si on veut une application naturelle de la surface recollée dans $X$, qui sur $\Delta_D$, $\Delta_{D'}$ se réduise (semble-t-il) aux plongements canoniques, il faut que $b$ et $\alpha$ se réalisent \add{ensemble} \uncertain{par} le \ill{} segment plongé dans $X$, ce qui imposerait $\alpha = a$ ! Mais si on convient de remplacer $\Delta_D$ par un disque un peu plus petit $\Delta^{*}_D$, \add{dont le bord} passe par $s$ (cf gribouillis ci-contre), et d'y \ill{} \ill{} \ill{} \ill{} $b$ en $b^{*}$, et faire le recollement le long de $b^{*}$ plutôt que $b$, et pareil pour le disque $\Delta_{D'}$ (de façon à faire coïncider les images de $a$ et $a'$ \struck{\ill{}} \struck{\ill{}} dans $X$), on voit qu'on arrive à \struck{\ill{} \ill{}} une application de la surface recollée dans $X$, même en choisissant $\alpha = a$.
\note{figure au crayon, dans la marge de gauche : un disque traversé de deux cordes qui se croisent en $s$ près du bord, un arc épais brun marqué $b^{*}$, légende $\Delta^{*}_D$ --- le « gribouillis » auquel il renvoie}

Le choix le plus naturel quand même (dans l'optique d'une application dans $X$) semblerait $\alpha = a'$. Donc, sur le terrain, cela donne une \emph{\uncertain{continuation}} d'une rive par la \uncertain{suivante} --- l'application $\mathcal{X} \to X$ est une \emph{immersion}.
\note{figure : deux pseudo-droites orange se croisant sur une droite vert foncé marquée $s$ (?) à son extrémité ; au-dessous, une rangée de points noirs, hachurée de vert clair au-dessus (marqué $\omega$) et de vert foncé au-dessous (marqué $\omega'$) ; plus bas, une droite vert clair fléchée marquée $\omega_D$. « continuation » et « immersion » sont soulignés}

Dans cette optique\supplied{, donc}, la rive de $D'$ i.e.\ l'arc sur $\partial\Delta_{D'}$ qui exprime l'opération de \struck{spécialisation} \add{générisation} vers $D$ est $a'$ --- i.e.\ celle qui « fait face » : la rive $a$, et non $a$ que on eût déduit par « transport parallèle » \ill{} \ill{} de la déformation de spécialisation $D \to D'$. Ce choix semble parfait --- mais on trouve à l'inverse,
\note{« générisation » est écrit au-dessus de « spécialisation » biffé. « et non $a$ » : la lettre est \uncertain{peu nette} et devrait désigner l'autre rive ; nous transcrivons ce que la page semble porter}

\page{66}
c'est qu'il semble que sur l'autre $\mathcal{X}$ obtenue par recollement, on n'ait aucune interprétation des « trous » (ou sommets, quand on a bouché les trous) en termes de la géométrie de $\underline{\Sigma}$.

\struck{\ill{}} Il faudrait quand même que je \uncertain{revoie} \ill{} \ill{}, pour voir si je ne m'étais pas \uncertain{mélangé} les pinceaux \ill{} \ill{} \ill{} (pas en faisant le recollement), et qui \ill{} \ill{} continuent une orientation choisie $\omega$ de $\Delta_D$ (défini par une or.\ $\omega_D$ de $D$) en passant de $\Delta_D$ dans $\Delta_{D'}$ via l'arc de transition, on trouve \struck{que} une orientation $\omega'$ de $\Delta_{D'}$, qui induit sur $D'$ une orientation $\omega_{D'}$ qui est \struck{\ill{}} \add{\ill{}} de celle déduite de $\omega_D$ par \struck{\ill{}} transport parallèle ! Voilà donc \ill{} transport \ill{} \ill{}, \ill{} vérifié \ill{} \ill{} \ill{} dans les deux \uncertain{sens} !
\note{le mot biffé après « qui est » est barré de deux traits ; ce qui est écrit au-dessus ne se lit pas. La fin de la colonne de gauche se lit mal}

Venons-en à l'autre choix qui donne lieu à $\mathcal{X} \to X$, savoir $\alpha = a$. On est obligé ici d'ajuster comme j'avais dit $(\Delta_D, b)$ et $(\Delta_{D'}, a)$, pour faire se rencontrer $b$ et $a$ dans $b^{*} \cup a^{*} \subset X$. On trouve la figure
\note{figure : deux pseudo-droites orange se croisant sur une ligne de points noirs marquée $\delta$, un segment épais brun au point de croisement ; au-dessus, des hachures alternées vert clair et vert foncé, avec les orientations $\omega$ et « $\omega' = -\omega$ » en petites flèches courbes ; au-dessous, une courbe vert foncé en cloche fléchée à gauche ($\omega_{D'}$) et une droite vert clair fléchée à droite ($\omega_D$)}

i.e.\ on recolle \ill{} \ill{} --- à la \emph{même rive} le long de la ligne critique de spécialisation $\delta$, l'identification $b^{*} \to a^{*}$, comme tantôt, et l'iso.\ de transport. Ici, la combinaison (\add{d'objets} choix de structure polygonale sur l'arc des transitions associés à une p.r.) est opposée au choix \ill{}

\page{67}
de tantôt, pour \uncertain{réaliser} une op.\ de générisation $D' \to D$ par un « arc de \struck{spécialisation} générisation » sur $\partial\Delta_{D'}$ ; on prend maintenant, non la rive qui « fait face » : celle de $b$, mais par contre celle qui en résulte par transport parallèle. \struck{Je \ill{}} \add{\uncertain{Quand} à l'isom.}\ de transition $\underline{\omega}_{\Delta_D} \to \underline{\omega}_{\Delta_{D'}}$ pour les orientations, \uncertain{faisant} \ill{} \ill{} \ill{}
\[
\omega \text{ or.\ de } \Delta_D \longleftrightarrow \omega_D \text{ or.\ de } D
\]
\[
\omega' \text{ or.\ de } \Delta_{D'} \longleftrightarrow \omega_{D'} \text{ or.\ de } D',
\]
\note{une flèche coudée, légendée « transport », va de $\omega$ à $\omega'$}
on trouve \ill{} \ill{} que $\omega_{D'}$ est \emph{\uncertain{opposée}} à l'orientation \struck{de $D$} déduite de $\omega_D$ par transport parallèle !

Pendant ces deux dernières semaines, c'est \ill{} cette opération de recollage que j'ai \ill{}. Avant, j'avais pris \struck{\ill{}} \add{\ill{}} iso.\ de transition pour les orientations, le transport parallèle \struck{\ill{}} et comme \struck{\ill{}} recollage, si je me rappelle bien, le recollage

1°) de la \ill{} --- je croyais avoir interprété contradictoirement la construction de recollement de 1°), mais je m'étais mélangé les pinceaux \emph{deux} fois, il me semble :

1) je m'étais trompé sur les isom.\ de transition pour les orientations, et

2) je m'étais trompé également sur l'identification associée à une op.\ de générisation \add{pour $D'$} par un arc sur $\partial\Delta_{D'}$ ; on prenait au lieu de la rive « \uncertain{d'en face} » celui \uncertain{qui} \struck{\ill{}} se déduit de $a$ par transport parallèle, \uncertain{alors} que c'est l'autre \struck{qu'il} (l'antipodique), qu'il fallait prendre).

Pour résumer les quatre possibilités sont indiquées dans la liste suivante
\note{« opposée » est souligné. La colonne de droite reprend en « 1°) » la fin de la colonne de gauche ; les deux erreurs sont numérotées 1) et 2), le premier chiffre \uncertain{peu net}. La liste annoncée suit à la page 68}

\page{68}
A) Recollements donnant lieu à une application naturelle
\[
\mathcal{X} \to X,
\]
ce qui revient à \ill{}, ou l'isom.\ $\alpha : a \to \beta$ \add{(correspondant à l'antipodisme)} est obtenu par transport parallèle --- ou encore, combinatoirement, où l'isom.\ de transition $\omega_D \to \omega_{D'}$ est opposé au transport parallèle. Il y a deux possibilités (car tjrs $\beta = b$)

1) Cas $\alpha = a'$ --- on recolle \struck{l'arc de \ill{}} \struck{\ill{} \ill{} $b$ \ill{}} l'arc de spéc.\ $b$ sur $\Delta_D$ à un \add{l'arc} \struck{\ill{}} de \ill{} $a'$ sur $\Delta_{D'}$ qui est sur la rive \emph{\uncertain{opposée}} \struck{\uncertain{face}} à celle de $b$, i.e.\ opposé à celle obtenue par transport parallèle : \struck{C'est} Alors $\mathcal{X} \to X$ est une immersion. \struck{On repère}

Pour définir la structure polygonale sur l'\uncertain{ens.}\ des transitions \uncertain{partant} d'une p.r.\ donnée $D'$, on repère \struck{d'une face} une générisation par la rive \uncertain{de} celle vers laquelle on va, qui est \struck{dans la direction} : le trait bleu qui déplace $D'$ par générisation marque la rive « tirée » sur $D'$ pour la déplacer vers lui.
\note{le début de l'alinéa porte un signe biffé ; plusieurs mots sont récrits sur des biffures}

2) Cas $\alpha = a$, --- on recolle $b$ avec l'arc $a$ qui s'en déduit par transport parallèle. Maintenant $\mathcal{X} \to X$ est une immersion sur chaque face $F^{*}$, mais c'est un « retournement » le long d'un \uncertain{arc de} transition commun à deux faces.

On repère une op.\ de générisation par \struck{\ill{}} l'arc de gén.\ qui est opposé à la rive dans laquelle on \struck{\uncertain{pousse}} va --- le trait bleu marquant cette rive « poussée » sur $D'$ par la générisation, et la direction \uncertain{dans} \uncertain{laquelle} \ill{} le sommet envisagé.
\note{« retournement » : lecture \uncertain{probable} d'un mot rapide ; la fin de la colonne de droite se lit mal}

\page{69}
B) Recollements \add{ne} donnant pas lieu à une application naturelle $\mathcal{X} \to X$, i.e.\ l'isom.\ $\alpha : a \to \beta$ est \emph{opposé} à celui obtenu par transport par la composition de \ill{} \ill{} antipodisme --- ou encore, l'isom.\ de transition $\underline{\omega}(D) \to \underline{\omega}(D')$ est \struck{\uncertain{opposé}} celui du transport parallèle. Deux cas à envisager encore

1) Cas $\alpha = a'$. C'est le cas où on repère une op.\ de générisation \add{sur $D'$} par l'arc qui fait face à l'arc de spécialisation correspondant sur $D$ --- c'était un cas que je crois avoir un peu vite j'avais écarté (c'était vraiment un \uncertain{bâton-dans-les-roues}, rétrospectivement !).

2) Cas $\alpha = a$ --- c'est le cas \add{\uncertain{déjà} \uncertain{rejeté}} où on repère une op.\ de générisation par l'arc \supplied{(}déduit de l'arc de spécialisation correspondant par transport parallèle. Je me suis \uncertain{amplement} convaincu qu'il est impossible de « \uncertain{tourner} » une p.r.\ \ill{} \ill{} \uncertain{quelques} \ill{} \uncertain{successives} de l'opération \ill{} de « tourner autour d'un trou (ou sommet \ill{}) ».

Je crois me rappeler maintenant que j'avais écarté aussi le cas 1), car il donnait lieu \uncertain{aussi} à des circuits qui \uncertain{menaient} \ill{} \ill{} \ill{} \ill{} \ill{} d'une p.r.\ proche d'une \uncertain{supersingulière}, $D_i$, on s'écarte de $D_i$ au lieu de tourner autour comme on devrait normalement.
\note{« B) » est écrit sur un chiffre surchargé ; « ne » est ajouté au-dessus de la ligne. « bâton-dans-les-roues » est une lecture \uncertain{hasardeuse} d'un mot composé rapide. La deuxième colonne continue la dernière phrase de la première}

\struck{Si je viens de \uncertain{vérifier}, cet ennui se présente \uncertain{aussi} bien pour B1) que pour B2), et \ill{} pour A1), \ill{}, je viens de voir, pour A2 \uncertain{visiblement}, dans un \ill{}}
\note{ce dernier alinéa est barré de quatre traits obliques ; les renvois B1), B2), A1), A2 sont écrits sur des lettres surchargées}

\page{70}
\struck{cas en question, c'est dû au fait que je n'ai pas tenu compte des arcs de spécialisation exceptionnels (vers les p.r.\ supersingulières). Le cas, il y a deux variantes, suivant qu'on se permet de se spécialiser vers une pos.\ rel.\ supersingulière, ou qu'on l'évite, en passant directement à la position voisine \uncertain{opposée}}
\note{les neuf premières lignes sont barrées de trois traits obliques}

Je viens de vérifier, sauf erreur (j'ai fait \ill{} une d'erreurs encore dans le cours de la vérification !) \uncertain{que} les trois possibilités $A_1$, $A_2$, $B_1$ donnent lieu aux bons circuits autour d'une position supersingulière, il n'y a que $B_2$ que \uncertain{dérange} complètement (\uncertain{or} il semblerait que c'était \uncertain{sur} $B_2$ que je m'étais \uncertain{arrêté} pendant une \uncertain{paire} de semaines, Dieu sait pourquoi !). La plus intéressante, après $A_2$, semble bien être $A_1$, mais je viens encore d'essayer de comprendre les circuits auxquels il donne naissance, \uncertain{sans} y réussir.
\note{« $B_1$ » : l'indice, sous une surcharge, est \uncertain{peu net}. La colonne de droite de la feuille est blanche}

\page{71}
Je voudrais revenir à la question de faire une surface avec les p.r.\ \emph{stables}. Pour une telle $D$, $\Delta_D$ n'a comme arcs de transition que des arcs de \emph{spécialisation}. Comment recoller $\Delta_D$ et $\Delta_{D'}$, quand $D$, $D'$ sont déduits l'une de l'autre par balayage au-dessus d'un sommet ? Ici, il n'y a pas d'ambiguïté possible sur les arcs le long desquels faire le recollement, ce sont les antipodiques des arcs de spécialisation --- et c'est plus \uncertain{utile} (il faudrait les voir pour \uncertain{l'utiliser} pour recoller !). Ici, il est très tentant de remplacer encore $\Delta_D$, $\Delta_{D'}$ par $\Delta^{\wedge}_D$, $\Delta^{\ast}_{D'}$ déduits en balayant le bord jusqu'aux points $s$ « critiques » qui nous intéressent, \struck{\ill{}} \add{\ill{}} de recoller.
\note{« stables » et « spécialisation » sont soulignés ; un double trait et un point d'exclamation, en marge, marquent la parenthèse. Les exposants de $\Delta^{\wedge}_D$ et $\Delta^{\ast}_{D'}$ sont deux petits signes \uncertain{peu nets}. Figure à gauche : deux pseudo-droites, $D'$ (vert foncé) et $D$ (vert clair), qui se séparent en passant de part et d'autre d'un point où concourent trois pseudo-droites orange}

Prenant une figure symétrique (\uncertain{à} partir de la position semi-stable intermédiaire) on a à priori $8$ possibilités de recollement, suivant \struck{\ill{}} les quatre possibilités d'échanges dans $\lbrace a, a' \rbrace \times \lbrace b, b' \rbrace$, puis le choix d'orientation pour recoller $\alpha \in \lbrace a, a' \rbrace$ avec $\beta \in \lbrace b, b' \rbrace$. Mais déjà le choix d'un couple $(a, b')$ ou $(b, a')$ est exclu, à cause de la \uncertain{nécessaire} symétrie à respecter. Donc il reste deux couples $(a, b)$, $(a', b')$, et pour chacun choix parmi les deux possibilités de recollage --- transport parallèle, ou l'opposé. Mais \uncertain{moyennant} \struck{application} transformation, par antipodisme \uncertain{des} \ill{} \ill{} droite (qui ne change \ill{} la surface $\mathcal{X}$ obtenue par recollement \struck{\ill{}} des $D \in \Sigma$ \uncertain{disposition} de $L$ \uncertain{dedans}) $(a', b')$ \uncertain{équivaut} à $(a, b)$.
\note{figures, en haut de la colonne de droite : un disque au bord vert clair traversé de deux pseudo-droites rouges qui se croisent près du haut et du bas, des arcs épais noirs marqués $a'$ et $a$ ; une droite au crayon le traverse et touche un second disque esquissé au crayon. Au-dessous, un disque au bord vert foncé, deux pseudo-droites rouges se croisant en haut et en bas, entourées de petits arcs vert foncé et vert clair, avec des arcs épais noirs marqués $b^{\ast}$, $a'$ en haut, $b'$, $b$ et $a$ en bas. La dernière ligne de la page est écrite sous une ligne raturée}

\page{72}
$N = 4n$
\[
\sigma_0 \sigma_1 \cdots \sigma_{N-1} \quad \Big| \quad \sigma_N = \sigma_0
\]
\[
\rho_1 = \sigma_0\sigma_1 \quad \rho_2 = \sigma_1\sigma_2 \quad \cdots \quad \rho_N = \sigma_{N-1}\sigma_N = \sigma_{N-1}\sigma_0
\]
\[
\rho_1\rho_2 = \sigma_0\sigma_2 \quad \rho_2\rho_3 =
\]
$\sigma_i^2 = 1$, $\rho_i^2 = 1$ i.e.\ deux $\sigma_i$ consécutifs (circulairement) commutent
\note{notes au crayon sur un listing d'ordinateur, la feuille tenue à l'envers du listing ; « $\rho_2\rho_3 =$ » reste sans second membre. Croquis au crayon : un polygone régulier (octogone) rayonné depuis son centre, une ligne brisée en zigzag épaissie le long d'une partie de son bord et des faisceaux de droites autour ; plusieurs faisceaux de droites sur une droite, avec les légendes « $n-1$ » et « $(n-1) - (N-1) - (N'-1)$ » \uncertain{lues en partie} ; une pseudo-droite fusiforme esquissée en haut de la page}

\page{73}
\note{figures en tête de la page : un réseau de hachures vert foncé et jaune coupé par deux pseudo-droites rouges qui se croisent sur une droite pointillée, un petit segment épais noir au croisement ; légendes $\omega_{D'}$, $\omega_D$ (flèches), $\omega_F$ et $\omega_{F'}$ (flèches courbes). Une accolade renvoie à une seconde figure analogue, légendée « ou $(a', b')$ » et, à droite, « \uncertain{Cas} un peu vicieux, et sans intérêt ! »}

\struck{\ill{}} Cas $(a, b)$ \struck{\ill{}} de prolongation des rives.

1) On choisit comme recollement le transport parallèle. On trouve un $\mathcal{X}$ \struck{avec} $\varphi : \mathcal{X} \to X$ qui est une immersion. Dans $\mathcal{X}$, on a un graphe top.\ $K$, et les faces de $(\mathcal{X}, K)$ correspondent aux p.r.\ stables. \struck{\ill{}} On a un $L = \varphi^{-1}(X_1)$, qui matérialise la corr.\ entre faces et p.r.\ stables --- en remplaçant chaque face $F$ par la face $F^{\sim}$ obtenue en la « renflant » aux pts qui « critiquent » de $S_K$ i.e.\ les pts de $K \cap S(L)$.
\note{« de prolongation des rives » est souligné ; le numéro qui précède « Cas » est surchargé, \uncertain{peut-être} un 2. $S_K$ : lecture \uncertain{probable}}

\marginal{NB Ici, on \ill{} $\Delta^{\ast}_D \cap$ \ill{} \ill{} $\Delta_{D'}$ \ill{} $D \cap \Delta_{D'}$}
\note{note écrite en biais entre les deux colonnes ; elle se lit mal}

Partons d'une orientation $\omega = \omega_F$ de $F = \Delta^{\ast}_D$, ou d'une $\omega_D$ orientation de $D$, prenons l'orientation $\omega'$ de $\Delta^{\ast}_{D'}$ déduite par transport via le recollement, d'où orientation $\omega_{D'}$ de $D'$, on constate avec \uncertain{surprise} que c'est \emph{l'opposée} de celle obtenue par transport parallèle !
\marginal{L'ennui, c'est que les circuits \uncertain{acceptables} \uncertain{ne} s'obtiennent pas pour les \uncertain{bons} circuits supersinguliers !}

2) Recollement par l'iso.\ \add{$a \simeq b$} opposé au transport parallèle. Cette fois, c'est au niveau de la transition $\underline{\omega}_F \to \underline{\omega}_{F'}$, interprétée comme $\underline{\omega}_D \to \underline{\omega}_{D'}$, qu'on obtient le transport parallèle. Contrairement à ce que je \add{croyais} (\uncertain{voyais}), on trouve les bons circuits autour des positions supersingulières.
\note{« l'opposée » est souligné ; la note marginale, à droite, est séparée du texte par un trait. « croyais » est écrit au-dessus de « (voyais) », qui n'est pas biffé}

\page{74}
Mais les circuits plus généraux restent mystérieux pour moi --- je ne sais s'ils sont \uncertain{valables} !

Les deux procédés correspondent à la \uncertain{même} structure polygonale sur l'ens.\ des arcs de spécialisation \add{($=$ transition)} pour une p.r.\ stable $D$ --- mais il y a une autre structure polygonale, dans le fait que \struck{les \ill{}} les arcs de spécialisation \emph{sur $D$} (au lieu de $\widetilde{D}$) sont mutuellement disjoints, donc on peut prendre la structure polygonale correspondante, dans la réalisation top.\ \uncertain{usuelle} de $D$ \uncertain{linéaire}. L'ennui, c'est que cette fois-ci, un polygone \struck{\ill{}} \add{d'ordre $n$} \uncertain{n'est} \uncertain{pas} \uncertain{réalisé} géométriquement \uncertain{comme} un $2n$, \struck{\ill{}} \uncertain{mais} \ill{} correspondant \uncertain{plutôt} par un disque $n$-gonal \uncertain{plus} \uncertain{ou} \uncertain{moins} clair. Mais peut-être est-ce une indication pour obtenir \emph{simultanément} des deux façons, \uncertain{antipodiques}, les disques $\Delta^{\ast}_D$, $\Delta^{\ast}_{D'}$, de façon à obtenir un \uncertain{ruban} de Möbius peut-être ?
\note{« sur $D$ », « simultanément » et « des deux façons » sont soulignés ; « d'ordre $n$ » est écrit au-dessus d'un mot biffé. Le passage entre « un polygone » et « Mais peut-être » se lit mal}

L'idée maintenant c'est de prendre le découpage de $X$ suivant la partie commune de $D$ et $D'$ (cf.\ figure) se donne \uncertain{dans} un ruban de Möbius, \struck{\ill{} \ill{}} \add{\uncertain{divisé} \uncertain{par} \uncertain{la} ligne} pointillée en une bande rectangulaire. Un \uncertain{prend} les \uncertain{orientations} \add{(\uncertain{twistées} par $\omega(D)$ ?)} des rubans de Möbius, \add{\uncertain{des} \uncertain{orientations} de $D'$} (\uncertain{ordre} $2$ des rubans) qui est formé de deux rectangles recollés en $2$ paires d'arêtes. Sauf erreur, c'est l'identification \struck{\ill{}} simultanée de $\hat{a}$ avec $b$, \struck{\ill{}} de $\hat{a}'$ avec $b'$, suivant le transport parallèle. Mais ce qui est \uncertain{décisif}, en faisant tous ces recollements pour \uncertain{chaque} \ill{}, \uncertain{au} lieu de $\Delta^{\ast}_D$, c'est celle bien de la surface cellulaire obtenue par le graphe $\Gamma_{\Sigma}$ et sa structure de polygonalisation et de transport ?, obtenue
\note{figure en tête de la colonne de droite : deux pseudo-droites, vert foncé et vert clair, confondues aux deux bouts et s'écartant en fuseau au milieu, une ligne pointillée noire le long de l'axe du fuseau, deux pseudo-droites rouges se croisant sur cet axe. La colonne se lit mal par endroits ; les lettres $\hat{a}$, $\hat{a}'$ portent un chapeau ajouté et sont écrites sur un signe biffé}

\page{75}
par recollement \struck{de} polygones \emph{d'ordre $n$}, non $2n$ ? \emph{Celle}-là \uncertain{première} ne serait-elle pas plutôt un revêtement d'ordre deux de la seconde, ramifié en chacun des pts des centres --- en effet, on voit \struck{\uncertain{difficilement}} \add{\uncertain{que} \uncertain{les}} antipodismes dans les $\Delta^{\ast}_D$ sont compatibles avec les identifications, et \uncertain{on} \uncertain{peut} \uncertain{passer} \uncertain{au} quotient, \uncertain{cela} \uncertain{va} donner \uncertain{justement} la carte cellulaire envisagée\ldots\ On peut aussi, \struck{\ill{}} \uncertain{plutôt} \struck{\ill{}} interpréter la nouvelle carte comme associée à un graphe \add{(\uncertain{défini} cellulairement)} dont les sommets sont les p.r.\ stables, que pour deux p.r.\ $D$, $D'$ adjacentes, l'on va maintenant envisager \emph{deux} flèches ou arêtes entre eux, savoir $a_D : D \to D'$ (de \uncertain{l'arête} $a$ vers $b'$) et $a_{D'} : D' \to D$ (de l'arête $b$ vers $a'$), munies chacune d'une orientation --- le balayage de $D$ vers $D'$, et de $D'$ vers $D$, qu'on \uncertain{identifie}
\note{« d'ordre $n$ » et « Celle » sont soulignés ; « deux » est souligné. Plusieurs mots de cette colonne, rapide, sont d'une lecture incertaine}

Maintenant, les arcs de transition (\uncertain{ou} \uncertain{plutôt}) pour $D$ i.e.\ sur $\partial\Delta_D$, sont les arcs de spécialisation \emph{et leurs antipodiques}, organisés en polygone de la façon habituelle. Une arête \uncertain{orientée} qui \uncertain{part} de l'un ou l'autre, en \uncertain{dit} \ill{} vers une position rel.\ adjacente.

La surface $\mathcal{X}$ obtenue, immergée dans $X$ par une immersion, et l'image inverse \add{$L$} de $X_1$ \struck{\ill{}} donne lieu à \struck{une} l'\uncertain{écriture} \uncertain{précédente} habituelle des \ill{} p.r.\ \uncertain{stables} associée à une face de $(\mathcal{X}, K)$.
\note{« et leurs antipodiques » est souligné}

Description combinatoire

1) graphe $\Gamma$ \quad $\left\lbrace \begin{array}{l} \text{sommets : p.r.\ stables } D \\ \text{arêtes orientées : paires } (D, a),\ a \text{ un arc de transition sur } D \\ \quad \text{(de spécialisation, ou antipodique d'un tel)} \end{array} \right.$
\note{sous « arêtes orientées », un « or. » est biffé ; la parenthèse est ajoutée au-dessous de la ligne}

\page{76}
\struck{antipodique d'un} opposée d'une arête orientée $(D, a)$ : l'arête $(D', a')$, où $D'$ est déduite de $D$ par balayage \add{\uncertain{via} \uncertain{l'arc}} \struck{\ill{}} $\operatorname{antip}(a)$, et $a'$ \add{antipodique de celle} déduite de $a$ par transport parallèle, i.e.\ « opposé à $a$ ». Origine de $(D, a)$ est $D$.
\note{« antipodique de celle » est ajouté au-dessous de la ligne ; la suite de la liste commencée page 75}

2) Structure polygonale sur l'ens.\ des arêtes orientées issues de $D$ : évidente.

3) Transport d'orientations : \emph{\uncertain{inverse}} du transport parallèle des orientations des p.r.\ !
\note{le mot souligné de 3) est écrit sur un autre ; lecture \uncertain{douteuse}}

Hélas on ne trouve pas les bons circuits supersinguliers --- ce n'est pas \uncertain{l'aboutissement}, \uncertain{si} alléchant que cela paraisse ! Si cependant on modifie comme donné de transport 3), en prenant pour les orientations le transport parallèle, on trouve les bons circuits supersinguliers (mais pas d'application naturelle $\mathcal{X} \to X$ !).

Il y a une autre (au moins !) possibilité de construction encore, c'est de recoller $a$ à $b'$, $b$ à $a'$, ainsi
\note{figure en tête de la colonne de droite : un réseau de hachures vert foncé et jaune coupé par deux pseudo-droites rouges qui se croisent sur une droite pointillée ; légendes $\omega_{D'}$, $\omega_{F'}$, $\omega_F$ (flèches), et au bas une droite vert clair fléchée $\omega_D$}

en utilisant pour le recollement le transport parallèle, de façon \uncertain{donc} à avoir \uncertain{toujours} $\mathcal{X} \to X$, mais cette fois-ci pas une immersion, mais \struck{des} application \uncertain{avec} \uncertain{pliure} le long des arêtes de $K$.

La description combinatoire est analogue à celle qui précède, la seule modification est dans 1) la description de l'opposée d'une arête, maintenant \add{\uncertain{c'est}} l'opposée $(D', a')$ de $(D, a)$, où $a'$ déduit de $a$ par transport parallèle. Pour 2), 3), \uncertain{on} \uncertain{retrouve}
\note{« 2), 3) » : le chiffre $3$ est écrit sur un autre, lecture \uncertain{probable}. La phrase se poursuit à la page 77}

\page{77}
Il s'ensuit que la qualité de $\mathcal{X}$ par antipodisme est la même que précédemment, et que les suites des p.r.\ intervenant dans un circuit sont les \add{mêmes} \struck{\ill{}} pour les deux --- donc cette variante n'est pas plus satisfaisante, en ce qui concerne les circuits supersinguliers ! Il faudrait modifier 3), en prenant la prescription opposée pour définir le transport d'orientation, pour trouver \struck{des} bons circuits supersinguliers.
\note{« mêmes » est écrit au-dessous d'un mot biffé ; « qualité » est d'une lecture \uncertain{douteuse}}

Que se passe-t-il dans le cas où $\underline{\Sigma}$ est un syst.\ de \emph{droites} dans un \add{vrai} plan projectif réel $X$ --- quand on regarde le système des droites dans $X^{\vee}$ correspondant aux sommets de $\underline{\Sigma}$, i.e.\ les droites joignant deux à deux les pts $\delta_i$ correspondant aux droites $D_i$ ? On suppose que les $\delta_i$ \add{\uncertain{choisis}} sont non alignés $3 : 3$ (i.e.\ $\underline{\Sigma}$ stable), et que les sommets \add{(distincts des $\delta_i$)} (de ce syst.\ de $\frac{n(n-1)}{2}$ droites) soient simples (quitte à perturber un peu le syst.\ des $\delta_i$). Donc il n'y a que des sommets d'ordre $2$, en dehors des sommets « supersinguliers » $\delta_i$.
\note{« droites » est souligné ; « vrai » est ajouté au-dessus de « plan projectif » ; le « $2$ » de « d'ordre $2$ » est d'une lecture \uncertain{probable}. La phrase passe d'une colonne à l'autre après « non »}

Considérons une face $F$ de $\check{\underline{\Sigma}}$, par définition \uncertain{dans} une p.r.\ $D = D_F$, et considérons les arêtes incidentes \add{$a_\alpha$ de $F$, $= \Delta_\alpha \cap F$}, qui correspondent à \emph{certaines} spécialisations $D \to D_\alpha$ (où $\alpha$ est l'indice d'une droite \add{$s_\alpha$} de $\Sigma$, $\Delta_\alpha$ la droite correspondante dans $X^{\vee}$). Le sommet $s_\alpha$ est le sommet d'un triangle de spécialisation pour $D$, dont $a_\alpha$ l'arc de spécialisation correspondant. Je voudrais
\note{« certaines » est souligné. $\check{\underline{\Sigma}}$ rend son $\Sigma$ doublement souligné surmonté d'un accent inversé (le système dual). La phrase se poursuit à la page 78}

\page{78}
interpréter la structure polygonale sur l'ens.\ des $a_\alpha$, en termes des $s_\alpha$ correspondants. Notons déjà qu'il y a correspondance entre les \struck{structures} \add{orientations} de $F$, et celles de $D$, et que le transport continu de l'orientation de $F$ vers une face adjacente $F'$ (le long de $a_\alpha$ disons) correspond au transport continu (ou « parallèle ») de l'orientation de $D$ au cours du balayage correspondant. \struck{(D.} Ce dessin paraît rendre plausible le choix « naïf » du transport d'orientation, qui dans le contexte des recollements de p.r.\ stables via balayage, \uncertain{du} \uncertain{point} de \uncertain{vue} semble en effet meilleur, du pt de vue des « circuits ». Le fait qu'on n'obtienne pas \uncertain{de} \uncertain{l'application} naturelle $\mathcal{X} \to X$ se rapproche du fait qu'il n'y a pas d'application naturelle $X^{\vee} \to X$.
\note{« orientations » est écrit au-dessus de « structures » biffé ; deux mots de l'avant-dernière phrase se lisent mal}

\note{figures dans la colonne de droite. En haut, au crayon : un arrangement de droites, dont un faisceau par un point, et une face $F$ orientée ($\omega_F$) ; sur son bord, les arêtes $a_\alpha$, $a_\beta$ et les sommets $\delta_0$, $\delta_1$, $\delta_2$, une flèche allant de $\delta_0$ à $\delta_1$. Au-dessous, en couleur : une droite vert clair $D = D_{\delta_0}$ orientée ($\omega_D$), des pseudo-droites vert clair en pointillé marquées $D_1$, $D_2$ au-dessus, et quatre paires de pseudo-droites rouges qui se croisent en des sommets dont $s_\alpha$, $s_\beta$ ; sur $D$, des segments épais noirs, l'un marqué $a'_\alpha$ ; une flèche courbe au point $x$ de $D$}

Partons avec une orientation $\omega_F$ de $F$ ($\mapsto \omega_D$ de $D$) et une arête $a_\alpha$ ($\mapsto$ \struck{$s_\alpha$} sommet de $\Sigma$) allant de $\delta_0 \in F$ ($\mapsto D = D_{\delta_0}$) au \add{\uncertain{point}} \uncertain{voisin} $\delta_1$ \ill{} droites \uncertain{vers} un pt $\delta_1$ voisin de $a_\alpha$ --- i.e.\ faisons pivoter un pt $x$ sur $D$ et une orientation de rotation en $x$, et faisons pivoter
\note{la fin de la colonne se lit mal ; la phrase se poursuit à la page 79. Les lettres grasses rendues $F$ sont un F doublé, écrit sur une lettre biffée}

\page{79}
$D = D_0 = D_{\delta_0}$ autour de $x$ dans le sens de l'orientation donnée, jusqu'à ce que $D_1$ passe par $s_\alpha$ (cette opération \uncertain{correspond}, en \uncertain{termes} de $\check{\underline{\Sigma}}$, à \uncertain{avancer} pour \uncertain{commencer} par un point $\delta_1$ voisin de $a_\alpha$). Puis suivant le long de $a_\alpha$ suivant l'orientation induite par $F$ --- ce qui \uncertain{signifie} pour $D_1$ pivoter autour de $s_\alpha$, dans un sens qu'il faudra expliciter \struck{\ill{}} (déterminé par l'orientation de $D_1$ déduite de celle de $D_0$ par transport, et la rive privilégiée dont on dispose. Notons que l'orientation \add{$\omega_\alpha$} de $\Delta_\alpha$ fixé par $a_\alpha$, celle \struck{\ill{}} \add{$\omega_{\delta_1}$} du point $\delta_1$ sur $\Delta_\alpha$, et la rive $\rho_\alpha$ de $\Delta_\alpha$ en $s$ déterminée par $F$, forment un syst.\ compatible au sens des \uncertain{flèches}, et qu'il s'ensuit que les orientations correspondantes au \uncertain{contact} du contact $(s_\alpha, D_1)$, $\omega = \omega_\alpha \in \underline{\omega}(X, s_\alpha)$, $\omega_{\delta_1} = \omega_{D_1}$, et $\rho'_\alpha = {}$rive de balayage, forment un système \emph{antiscompatible}, i.e.\ l'orientation $\omega_{D_1}$ de $D_1$ est \emph{opposée} à celle induite par $\omega$ sur la rive de balayage (cf.\ figure) \add{\uncertain{dit} \uncertain{autrement}} \uncertain{ceci} \uncertain{permet} de propager le long de $a_\alpha$ \ill{} \uncertain{faisant} tourner
\note{« antiscompatible » (lecture \uncertain{probable}) et « opposée » sont soulignés ; plusieurs mots de la fin de la colonne, écrits sur des ajouts interlinéaires, se lisent mal}

$D_1$ autour de $s_\alpha$ dans la direction indiquée, jusqu'à ce moment où, en traversant un premier sommet (qui sera $s_\beta$), qui est aussi celui où le point \uncertain{variable} $\delta$ sur $a_\alpha$ \uncertain{arrive} au sommet $\delta_2$, rencontre de \struck{\ill{}} $a_\alpha$ et $a_\beta$ i.e.\ de $\Delta_\alpha$ et $\Delta_\beta$ [Bien sûr, \struck{entre les \ill{} \ill{}} \struck{déplacements de $D_\delta$, \ill{}} spécialisations correspondant à des $a_\beta$, il peut y avoir, dans un sens comme dans les autres ($\omega_D$ ou l'orientation opposée) \uncertain{bien} d'autres arcs de spécialisation, mais correspondant à des sommets $s_\gamma$ qui ne sont pas atteints au cours du balayage.] À partir de $s_\beta$, $D_\beta$ on continue comme précédemment.
\note{« (qui sera $s_\beta$) » est ajouté au-dessous de la ligne. Les crochets sont de sa main}

\page{80}
\note{page de dessins en couleur, avec une seule légende : quatre disques au bord jaune ou vert, chacun traversé d'un arrangement de pseudo-droites orange (pleines ou pointillées) et d'un faisceau de lignes grises au crayon allant d'un point du bord au point antipodique. Sur le premier, une corde rouge épaisse suit le diamètre ; sur le deuxième et le troisième, deux arcs rouges épais forment un fuseau entre deux points antipodiques ; sur le quatrième, le fuseau s'est refermé en une ligne grise le long du diamètre. Des flèches mènent du premier aux trois autres, avec les légendes « $\frac12$ circuits autour de pts d'ordre $1$, $2$ » et « ou $0$ »}

\end{document}
